\begin{longtable}{rp{2.5cm}|p{6.5cm}p{5.5cm}}
    \hline
    \rowcolor{head} \multicolumn{1}{c}{\textbf{Nr.}} & \multicolumn{1}{c|}{\textbf{Testfall}} &
    \multicolumn{1}{c}{\textbf{Beschreibung}} & \multicolumn{1}{c}{\textbf{Erwartetes Ergebnis}} \\
    \hline
    1 & Consistent\linebreak \Fachbegriff{Repository} & Das \Fachbegriff{Repository} wird über definierte Schnittstellen zu verschiedenen Laufzeiten
    angesprochen. & Daten sind vorhanden und können nur vom \Fachbegriff{Repository} modifiziert werden. \\
    \hline
    \rowcolor{odd}2 & Task statistics validated & Es werden \Abkuerzung{SQL}-Abfragen bei verschiedenen Datenbankständen
    durchgeführt. & Korrekte Ergebnisse der Tasks-Statistik werden geliefert. \\
    \hline
    3 & Connectivity & Beim Aufrufen von \Abkuerzung{JSON} wird das Netzwerk-Traffic überwacht.  & Client-Server Verbindung besteht
    (\acs{JSON}HttpRequest auf "pending"). \\
    \hline
    \rowcolor{odd}4 & Incoming\linebreak event & Es wird ein Event publiziert. & Im Frontend wird die JavaScript-Methode
    \Methode{onMessage} angesprochen. \\
    \hline
    5               & Data transfer & Es wird ein Event mit Inhalt publiziert. & Event-Daten sind im \Abkuerzung{JSON}-Format. \\
    \hline
    \rowcolor{odd}6 & Frontend event handling & Es wird ein Event publiziert. & Task-Event sorgt für Aktualisierung der
    \Abkuerzung{UI}-Task-Status-Anzeige. \\
    \hline
    7 & User scoping & Es werden Events mit unterschiedlichen Zielgruppen/-personen publiziert. & Nur gewünschte Benutzer werden
    benachrichtigt. \\
    \hline
    \rowcolor{odd}8 & Accessibility & Die Weboberfläche wird S/W betrachtet.   & Status-Anzeigen sind barrierefrei Nutzbar.    \\
    \hline
    9 & Browser\linebreak compatibility & Die Weboberfläche wird mit den vier "major browser"\footnotemark{} getestet. & Das
    Design funktioniert bei allen angemessen.  \\
\end{longtable}